\chapter{Setup and definitions}

We work with $n \times n$ real matrices ($n \ge 2$), indexed by
$[n] = \{1, 2, \dots, n\}$ throughout. The five definitions below set up
both the matrix-side object of study (orthogonal upper Hessenberg
matrices), its combinatorial shadow (the support digraph), and the two
parameterisations (Givens factorisation, active set) that allow the
matrix and combinatorial sides to be compared.

\begin{definition}\label{def:OH}
  \lean{HessenbergDigraphs.OrthogonalHessenberg}
  \leanok
  A matrix $Q \in \mathbb{R}^{n\times n}$ is \emph{upper Hessenberg} if
  $Q_{ij} = 0$ whenever $j+1 < i$, and \emph{orthogonal} if
  $Q^\top Q = I$. We write $\mathcal H_n$ for the set of orthogonal
  upper Hessenberg matrices in $\mathbb{R}^{n\times n}$.
\end{definition}

\begin{definition}\label{def:supp}
  \lean{HessenbergDigraphs.OrthogonalHessenberg.supportDigraph}
  \leanok
  \uses{def:OH}
  The \emph{support digraph} of $Q$, denoted $D(Q)$, is the directed
  graph on $[n]$ with an arc $i \to j$ whenever $Q_{ij} \ne 0$. The
  classification problem is to describe, up to digraph isomorphism, the
  set $\{D(Q) : Q \in \mathcal H_n\}$.
\end{definition}

\begin{definition}\label{def:sign}
  \lean{Matrix.SignEquiv}
  \leanok
  Two matrices $Q, Q' \in \mathbb{R}^{n\times n}$ are \emph{sign-equivalent},
  written $Q \simeq_{\mathrm{sign}} Q'$, if $Q' = D_L Q D_R$ for some
  diagonal matrices $D_L, D_R$ with entries in $\{\pm 1\}$.
  Sign-equivalent matrices have identical zero / non-zero patterns, so
  $D(Q) = D(Q')$.
\end{definition}

\begin{definition}\label{def:givens}
  \lean{Matrix.givensRotation, Matrix.givensProduct}
  \leanok
  For $k \in \{0,\dots,n-2\}$ and $\theta \in \mathbb{R}$, the \emph{Givens
  rotation} $G_k(\theta) \in \mathbb{R}^{n\times n}$ is the identity with the
  $2\times 2$ block at rows / columns $(k, k+1)$ replaced by
  $\bigl[\begin{smallmatrix}\cos\theta & -\sin\theta\\ \sin\theta & \cos\theta\end{smallmatrix}\bigr]$.
  The \emph{ordered Givens product} of an angle vector
  $\theta = (\theta_0, \dots, \theta_{n-2})$ is
  $Q_n(\theta) := G_0(\theta_0) G_1(\theta_1) \cdots G_{n-2}(\theta_{n-2})$.
  An angle vector is called \emph{unreduced} if every $\sin\theta_k$
  is non-zero; we write $\Theta_n$ for the set of such vectors.
\end{definition}

\begin{definition}\label{def:active}
  \lean{HessenbergDigraphs.ActiveSet, HessenbergDigraphs.Arc}
  \leanok
  \uses{def:givens}
  For $\theta \in \Theta_n$, the \emph{active set} is
  $S(\theta) := \{ k \in [n-1] : \cos\theta_{k-1} \ne 0\}$.
  For any $S \subseteq [n-1]$, set
  $R(S) := \{1\} \cup \{k+1 : k \in S\}$ and $C(S) := S \cup \{n\}$,
  and define $D_n(S)$ on $[n]$ to have arcs $j+1 \to j$ for
  $1 \le j \le n-1$ (\emph{spine}) plus $i \to j$ whenever
  $i \in R(S)$, $j \in C(S)$, $i \le j$ (\emph{overlay}).
\end{definition}

\chapter{Main theorems}

The classification of support digraphs of unreduced real orthogonal
upper Hessenberg matrices is the conjunction of four formal statements:
a universality result on the matrix side, a bridge that identifies the
matrix support with the combinatorial digraph, a rigidity result on the
combinatorial side, and the apex theorem assembling the three.

\begin{theorem}[Universality]\label{thm:univ}
  \lean{HessenbergDigraphs.OrthogonalHessenberg.exists_givensFactorization}
  \leanok
  \uses{def:OH, def:givens, def:sign}
  Every unreduced $Q \in \mathcal H_n$ is sign-equivalent to a canonical
  Givens product: there exists $\theta \in \Theta_n$ with
  $Q \simeq_{\mathrm{sign}} Q_n(\theta)$.
\end{theorem}
\begin{proof}\proves{thm:univ}\leanok\uses{def:OH, def:givens, def:sign}
  Induction on $n$. Right-multiply $Q$ by the transpose of a Givens
  factor chosen to clear the last row to $e_n^\top$; orthogonality
  forces the last column to $e_n$ as well, so the principal
  $(n{-}1) \times (n{-}1)$ corner is itself orthogonal upper Hessenberg,
  and the induction hypothesis applies. The accumulated sign
  discrepancies along the inductive peel-off are absorbed into the
  diagonal $D_L, D_R$ of the sign-equivalence relation.
\end{proof}

Universality reduces every unreduced $Q$ to the canonical form
$Q_n(\theta)$ up to sign. The next step is to read the support pattern
of $Q_n(\theta)$ off the active set $S(\theta)$ in a purely
combinatorial way.

\begin{theorem}[Bridge]\label{thm:bridge}
  \lean{HessenbergDigraphs.digraph_model}
  \leanok
  \uses{def:supp, def:givens, def:active}
  For every $\theta \in \Theta_n$, the matrix support agrees with the
  combinatorial digraph:
  $D(Q_n(\theta)) = D_n(S(\theta))$.
\end{theorem}
\begin{proof}\proves{thm:bridge}\leanok\uses{def:supp, def:givens, def:active}
  Two inclusions, both stated entry-wise on $Q_n(\theta)_{ij}$. The
  subdiagonal entries are $\sin\theta_j \ne 0$ by the unreduced
  hypothesis, exhibiting all spine arcs. For an upper-triangular entry
  $i \le j$, a row / column-vanishing analysis on the Givens product
  shows that $Q_n(\theta)_{ij} \ne 0$ iff $\cos\theta_{i-1} \ne 0$ and
  $\cos\theta_j \ne 0$, which by definition of $S(\theta)$ is exactly
  the overlay condition $i \in R(S(\theta))$, $j \in C(S(\theta))$.
\end{proof}

The bridge transports the classification problem from matrices to
combinatorics: $D(Q)$ is determined by the active set $S \subseteq [n-1]$.
The combinatorial side is now governed by the rigidity of the
$D_n(S)$ family.

\begin{theorem}[Rigidity]\label{thm:rigid}
  \lean{HessenbergDigraphs.rigid}
  \leanok
  \uses{def:active}
  Let $S, S' \subseteq [n-1]$ with $|S|, |S'| \ge 2$. If
  $D_n(S) \cong D_n(S')$ as digraphs, then $S = S'$.
\end{theorem}
\begin{proof}\proves{thm:rigid}\leanok\uses{def:active}
  The key combinatorial fact is the \emph{cut lemma}: in $D_n(S)$, the
  only arc from $\{k{+}1, \dots, n\}$ to $\{1, \dots, k\}$ is the
  spine arc $k{+}1 \to k$. This forces every isomorphism to fix the
  unique directed Hamilton cycle
  $1 \to n \to n{-}1 \to \cdots \to 2 \to 1$. A pigeonhole argument on
  the cut, combined with degree analysis at the maximum-in-degree
  vertex (which is $n$ when $|S| \ge 2$), forces the isomorphism to
  fix $n$, after which a downward induction along the cycle forces the
  identity. The case $|S| = 1$ is the only failure: the cyclic rotation
  $\rho^t$ exhibits $D_n(\{t\}) \cong D_n(\{n - t\})$, with no further
  coincidences.
\end{proof}

The three results above plug together as follows.

\begin{theorem}[Classification]\label{thm:class}
  \lean{HessenbergDigraphs.classification}
  \leanok
  \uses{thm:univ, thm:bridge, thm:rigid}
  Fix $n \ge 2$. The classification of support digraphs of unreduced
  orthogonal upper Hessenberg matrices is the conjunction of
  realisability (every $S \subseteq [n-1]$ arises as $S(\theta)$),
  bridge ($D(Q_n(\theta)) = D_n(S(\theta))$), and rigidity
  ($|S|,|S'| \ge 2$ and $D_n(S) \cong D_n(S')$ imply $S = S'$).
\end{theorem}
\begin{proof}\proves{thm:class}\leanok\uses{thm:univ, thm:bridge, thm:rigid}
  Realisability is witnessed by the explicit angle choice
  $\theta_k = \pi/4$ if $k+1 \in S$, $\theta_k = \pi/2$ otherwise; this
  gives $\sin\theta_k \ne 0$ always and $\cos\theta_k \ne 0$ exactly on
  $S$. Bridge is Theorem~\ref{thm:bridge}, applied pointwise. Rigidity
  is Theorem~\ref{thm:rigid}. The Lean-side declaration is a one-line
  term-level packaging of these three components.
\end{proof}

\chapter{Enumeration consequences}

Combining rigidity with the singleton coincidence
$D_n(\{t\}) \cong D_n(\{n-t\})$ yields a closed-form count of
isomorphism classes. Both counts follow the same pattern: enumerate
labelled configurations (subsets of $[n-1]$, or its loopless
restriction to subsets of $\{2,\dots,n-2\}$ with no consecutive
elements), then quotient by the singleton coincidence using rigidity to
keep all $|S| \ge 2$ classes distinct.

\begin{theorem}[Connected count]\label{thm:count}
  \lean{HessenbergDigraphs.card_isoClasses_quot}
  \leanok
  \uses{thm:rigid}
  For $n \ge 2$, the number of non-isomorphic weakly connected support
  digraphs of $n \times n$ unreduced orthogonal upper Hessenberg
  matrices is $N_n = 2^{n-1} - \lfloor (n-1)/2 \rfloor$.
\end{theorem}
\begin{proof}\proves{thm:count}\leanok\uses{thm:rigid}
  Labelled configurations are subsets $S \subseteq [n-1]$, of which
  there are $2^{n-1}$. Subsets with $|S| \ge 2$ are rigid
  (Theorem~\ref{thm:rigid}), so each is its own isomorphism class. The
  empty subset is rigid alone. The $|S| = 1$ subsets pair up under
  $\{t, n-t\}$, contributing $\lceil(n-1)/2\rceil$ classes. The total
  collapses to $2^{n-1} - \lfloor (n-1)/2 \rfloor$.
\end{proof}

\begin{theorem}[Loopless count]\label{thm:count_loopless}
  \lean{HessenbergDigraphs.card_looplessClasses_quot}
  \leanok
  \uses{thm:rigid}
  For $n \ge 3$, the number of non-isomorphic loopless connected
  support digraphs is
  $N_n^{\mathrm{loopless}} = F_{n-1} - \lfloor (n-3)/2 \rfloor$,
  where $F_k$ is the $k$-th Fibonacci number.
\end{theorem}
\begin{proof}\proves{thm:count_loopless}\leanok\uses{thm:rigid}
  The loopless condition forces $R(S) \cap C(S) = \emptyset$, which
  translates to $S \subseteq \{2, \dots, n-2\}$ with no two consecutive
  elements. The number of such subsets is the Fibonacci number
  $F_{n-1}$ via the standard binary-strings bijection. The same
  rigidity-plus-singleton-pairing argument as in Theorem~\ref{thm:count}
  reduces this to $F_{n-1} - \lfloor (n-3)/2 \rfloor$ isomorphism
  classes.
\end{proof}
